\scanpage{53}
\mspage{46}

d'être compatible avec les structures $C,\mathrm{TC}$ (propriété d'exactitude \inserted{homotopique} à droite, \unsure{moralement|normalement}, alors que $u_*$ est exact à gauche et \deleted{\unsure{p}} en général pas du \deleted{tout} à droite), ni $u_!$ avec les structures $F,\mathrm{TF}$ (commentaire dual à donner, $u_!$ est exact à \uline{droite}, et n'a pas envie d'être ``homotopiquement exact à gauche'').

Démonstration de \unsure{(1) (2) (3)}. Les relations (1) sont évidentes par définition. Les relations \unsure{(1)|(2)}, (3) le sont à peine moins, vu la propriété \unsure{surtout|somme toute} générale suivante, dont la place serait au § 1 :

\uline{Proposition 2} Soit\deleted{\unsure{ent}}

\diagram{\begin{tikzcd}[ampersand replacement=\&]
C' \arrow[r,shift left,no head,"u_!"] \& \deleted{\unsure{C}}C\arrow[l,shift left,"u^*"]
\end{tikzcd}}

un couple de foncteurs adjoints $(u_!,u^*)$, et soit $\Phi'\subset\mathrm{Fl}\,C'$, $\Psi\subset\mathrm{Fl}(C)$. Alors

\[ (u_!(\Phi')\longleftrightarrow\Psi)\Longleftrightarrow(\Phi'\longleftrightarrow u^*(\Psi)) \]

\uline{Corollaire} (C'est le $=$ énoncé, avec seulement \unsure{des|aux|un} \unsure{notation|notations} changée\unsure{s}). Supposons que $u^*:C\to C'$ \unsure{admette} un adjoint à \uline{droite} $u_*$ (\unsure{pas} forcément un adjoint à gauche).

\diagram{\begin{tikzcd}[ampersand replacement=\&]
C'\arrow[r,shift right,no head,"u_*"'] \& C\arrow[l,shift right,"u^*"']{,}
\end{tikzcd}}

soit $\Phi\subset\mathrm{Fl}(C)$, $\Psi'\subset\mathrm{Fl}(C')$, alors

\[ (u^*(\Phi)\longleftrightarrow\Psi')\Longleftrightarrow(\Phi\longleftrightarrow u_*(\Psi')) \]

C'est la prop. 2 \deleted{appliquée \unsure{au}} \inserted{\unsure{via|vu|en}} \deleted{$\unsure{C},C^{\unsure{n}}$}

\diagram{\begin{tikzcd}[ampersand replacement=\&]
\deleted{C'_0}\arrow[r,shift left,"\deleted{u_{0!}}"] \& \deleted{C_0}\arrow[l,shift left,"\deleted{u_0^*}"]
\end{tikzcd}}

\deleted{\unsure{tous les objets}} avec $C'_0=C$, $C_0=C'$, $(u_{0!}:C\to C')\deleted{\unsure{=u}}=u^*$, et $(u_0^*:C'\unsure{-|\to}C)=u_*$.
% archive sheet number: 52
